\documentclass[10pt,a4paper]{amsart}
\usepackage[margin=19mm]{geometry}
\usepackage{amsmath,amssymb,mathtools}
\usepackage[scr=rsfs,cal=euler]{mathalfa}
\usepackage{xfrac}
\usepackage[unicode,hidelinks,bookmarks=false]{hyperref}
\newcommand{\ie}{i.e.\ }
\newcommand{\cf}{cf.\ }
\newcommand{\resp}{respectively\ }
\newcommand{\Subsection}{\subsection}
\providecommand{\nobreakdash}{\nobreak\hbox{-}}
\setlength{\emergencystretch}{2em}
\multlinegap0pt
\usepackage{bm}
\usepackage{bbm}
\usepackage{dsfont}
\usepackage{xspace}
\usepackage{cancel}
\usepackage{mathtools}
\usepackage{amsthm}
\makeatletter
\@ifundefined{theorem}{\newtheorem{theorem}{Theorem}}{}
\@ifundefined{lemma}{\newtheorem{lemma}[theorem]{Lemma}}{}
\@ifundefined{prop}{\newtheorem{prop}[theorem]{Proposition}}{}
\makeatother
\makeatletter
\newcommand{\N}{\mathbb{N}}
\newcommand{\R}{\mathbb{R}}
\newcommand{\Rm}{\operatorname{Rm}}
\expandafter\ifx\csname customthm\endcsname\relax
\newenvironment{customthm}[1]
  {\renewcommand\theinnercustomthm{#1}\innercustomthm}
\fi
\providecommand{\sol}{\operatorname{sol}}
\makeatother
\title[K\"ahler-Ricci Tangent Flows in the Analytic Minimal Model P]{K\"ahler-Ricci Tangent Flows in the Analytic Minimal Model Program}
\author{Longteng Chen, Max Hallgren, Lucas Lavoyer}
\date{Source: arXiv:2608.19152v1 (CC BY 4.0)}
\begin{document}
\maketitle

\noindent\emph{Source and licence.} K\"ahler-Ricci Tangent Flows in the Analytic Minimal Model Program. Version arXiv:2608.19152v1 (CC BY 4.0), distributed under the Creative Commons Attribution 4.0 International licence, as recorded in the arXiv licence metadata for this exact version and shown on the abstract page https://arxiv.org/abs/2608.19152v1. Full rights notice: licenses/arxiv-2608.19152-CC-BY-4.0.txt.\par\smallskip

\noindent\textit{Editorial scope.} Counted unit: the Prop whose source label is \texttt{prop:modelmetrics} in the article (source0.tex lines 1681--1709, source label \texttt{prop:modelmetrics}), delivered together with the article's own complete original proof of it (source0.tex lines 1710--1989, one \texttt{proof} environment of 280 lines). No mathematics is written, repaired or re-derived: statement and proof are the article's own text line for line. Required context retained uncounted: lem:basicACfacts (lines 735--755); lem:basicACfacts3 (lines 740--742); lem:eigenfunctiongrowth (lines 1313--1333); lem:eigenfunctiongrowth2 (lines 1324--1326); lem; equiv of potentials (lines 1620--1630); lem; equiv of potentials1 (lines 1622--1629); tech lemma on lie derivative (lines 1653--1662). Every definition, display and label of those lines is retained unchanged. Omitted from this package and not redistributed: 1--734, 743--1312, 1327--1619, 1630--1652, 1663--1680, 1990--3384. Nothing inside a retained excerpt is omitted or altered: every retained body is delivered exactly as the article's lines are written, so no omission receipt of any kind accompanies this package. Citations: the retained excerpts carry no citation command at all, so no bibliography entry of the article is transcribed and no reference is redistributed. Reference adaptation: none was needed and none was made; every reference inside the delivered unit resolves inside it, so no unresolved reference or citation is produced. Printed slips kept literal: no printed word, symbol, hypothesis or number of the counted statement or of its proof was altered. Layout and class: the article's own class is \texttt{amsart} with packages including amsmath, amsthm, verbatim, amscd, amssymb, setspace, enumitem, hyperref, exscale, color, todonotes, graphicx, tikz, mathrsfs; the wrapper is the lane's pinned \texttt{amsart} editorial wrapper at 10pt on A4, which restates the theorem-like environments the retained text needs and adds the article's own macro definitions that the retained text needs (\texttt{\textbackslash N}, \texttt{\textbackslash R}, \texttt{\textbackslash Rm}, \texttt{\textbackslash sol}), with extra wrapper packages (bm, bbm, dsfont, xspace, cancel, mathtools). The delivered numbering is the unit's own. Assets: no retained span contains a figure environment, a graphics-inclusion command, a PDF-page inclusion, a table or a TikZ picture; no image file is copied into this package, the package declares no asset dependency and no image file is redistributed with it. Licence: version arXiv:2608.19152v1 of this article is distributed under the Creative Commons Attribution 4.0 International licence, as recorded in the arXiv licence metadata for this exact version and shown on the abstract page https://arxiv.org/abs/2608.19152v1. A full rights notice is delivered at licenses/arxiv-2608.19152-CC-BY-4.0.txt. Characters: every retained excerpt is pure ASCII, so the delivered engine, XeLaTeX with its default Latin Modern text font, reports no missing character.
\begin{lemma} \label{lem:basicACfacts}
    \begin{enumerate}
        \item \label{lem:basicACfacts1} $b_{\sol,t}$ converge uniformly on $M_{\sol}$ as $t\nearrow 0$ to $\pi_{\sol}^{\ast}r$.

        \item \label{lem:basicACfacts2} For all $(x,t)\in M_{\sol} \times (-\infty,0)$, we have
        \begin{equation*}
            C(g_{\sol})^{-1}(b_{\sol,t}(x)+\sqrt{|t|})\leq b_{\sol,0}(x)+\sqrt{|t|} \leq (b_{\sol,t}(x)+\sqrt{|t|}).
        \end{equation*}

        \item \label{lem:basicACfacts3} For any $k\in \mathbb{N}$, there exists $C_k =C_k(g_{\sol})$ such that for all $(x,t)\in M_{\sol}\times (-\infty,0)$, we have
        \begin{equation*}
            |(\nabla^{g_{\sol,t}})^k\Rm(g_{\sol,t})|_{g_{\sol,t}}(x) \leq \frac{C_k}{(b_{\sol,t}(x)+\sqrt{|t|})^{k+2}}.
        \end{equation*}

        \item \label{lem:basicACfacts4}
        There exists $C(g_{\sol}) \in (1,\infty)$ such that for any $x\in M_{\sol}$ and $r\in (0,\infty)$,
        \begin{equation*}
            \operatorname{Vol}_{g_{\sol}}(B_{g_{\sol}}(x,r))\geq C(g_{\sol})^{-1}r^{2n}.
        \end{equation*}
    \end{enumerate}
\end{lemma}

\begin{lemma} \label{lem:eigenfunctiongrowth} Suppose that for some $\lambda \in [0,\infty)$, $v\in C^{\infty}(M_{\sol})\cap L^2(M_{\sol},\nu_{\sol})$ satisfies
\begin{equation} \label{eq:eigenfunction} \Delta_{f_{\sol}} v = -\lambda v .
\end{equation}
\begin{enumerate}
    \item If $(M_{\sol},g_{\sol})$ has bounded Ricci curvature and $\lambda \in [0,1]$, then there exists $C_0(g_{\sol}) \in (0,\infty)$ such that
    \label{lem:eigenfunctiongrowth0}
    \begin{equation*}
        |v|\leq C_0(g_{\sol})\left( \int_{M_{\sol}}v^2 d\nu_{\sol} \right)^{\frac{1}{2}}b_{\sol}^{2\lambda}.
    \end{equation*}

    \item If $(M_{\sol},g_{\sol})$ has bounded curvature, then there exists $C_0(g_{\sol},\lambda) \in (0,\infty)$ such that \label{lem:eigenfunctiongrowth1} 
    \begin{equation*} |v| +b_{\sol}^{-1}|\partial v|_{g_{\sol}} 
    \leq C_0(g_{\sol},\lambda) \left( \int_{M_{\sol}}v^2 d\nu_{\sol} \right)^{\frac{1}{2}}b_{\sol}^{2n-1+2\lambda}.
    \end{equation*}

    \item \label{lem:eigenfunctiongrowth2} If $\lambda \in [0,1]$ and $(M_{\sol},g_{\sol})$ is either AC or compact, then for each $k\in \mathbb{N}$, there exists $C_k(g_{\sol}) \in (0,\infty)$ such that the following hold:
    \begin{equation*}
    \sup_{M_{\sol}}\sum_{j=0}^k b_{\sol}^{j-2\lambda}|(\nabla^{g_{\sol}})^j v|_{g_{\sol}}\leq C_k(g_{\sol}) \left( \int_{M_{\sol}} v^2 d\nu_{\sol}\right)^{\frac{1}{2}}.
\end{equation*}
\end{enumerate}
\end{lemma}

\begin{lemma}\label{lem; equiv of potentials}
    There exists $C(g_{\sol}) \in (1,\infty)$ such that the following hold for all $h\in H$:
\begin{enumerate}
    \item \label{lem; equiv of potentials1}
    $e^{-C(g_{\sol})\|h\|_H} f_{\sol} \leq f_{h} \leq e^{C(g_{\sol})\|h\|_H}f_{\sol}$,
    \item \label{lem; equiv of potentials2} for all $x_0\in \arg\min f_{h}$, 
\begin{equation*}
    \frac{1}{2}(d_{g_{h}}(x_0,\cdot)-C(g_{\sol}))_+^2 \le f_{h}\le \frac{1}{2}(d_{g_{h}}(x_0,\cdot)+C(g_{\sol}))^2.
\end{equation*}
\end{enumerate}     
\end{lemma}

\begin{lemma}\label{tech lemma on lie derivative}
    Suppose $(Y_t)_{t\in [0,1]}$ is a smooth family of complete vector fields on a smooth manifold $M$, and for each $t\in [0,1]$, let $(\Phi_s^{Y_t})_{s\in\R}$ denote the flow of $Y_t$. For all $t\in [0,1]$ and $s\in \mathbb{R}$, define
    \begin{equation*}
        Z_s^t=\frac{\partial}{\partial t}\Phi_s^{Y_t}\circ \Phi_{-s}^{Y_t}.
    \end{equation*}
    Then we have
    \begin{equation*}
        Z_s^t=\int_0^s(\Phi_{s-\tau}^{Y_t})_{\ast}(\partial_t Y_t) d\tau.
    \end{equation*}
\end{lemma}

\begin{prop} \label{prop:modelmetrics} Assume that $(M_{\sol},g_{\sol})$ is AC or compact.

\begin{enumerate}

    \item \label{modelmetrics0} For all $h\in H$, 
\begin{equation*} 
   e^{-C(g_{\sol})\|h\|_H} \omega_{\sol} \leq  \omega_h \leq e^{C(g_{\sol})\|h\|_H}\omega_{\sol}
\end{equation*}

\item \label{modelmetrics1} For all $h,k\in H$ with $\|h\|_H,\|k\|_H \leq 1$, 
    \begin{equation*} |u_h-u_k|\le C(g_{\sol})\|h-k\|_H f_{\sol},\end{equation*}

\item \label{modelmetrics2.5} For any $\ell_1,\ell_2 \in \mathbb{N}$, $m\in \mathbb{N}$, and $\mu \in \mathbb{R}$, there exists $C_{\ell_1,\ell_2,m}(g_{\sol},\mu) \in (1,\infty)$ such that for any smooth tensor $T$ of type $(\ell_1,\ell_2)$ on $M_{\sol}$, any $s\in [0,1]$, and any $h\in H$ with $\|h\|_H \leq 1$, we have 
\begin{equation*}
    \sum_{j=0}^m \sup_{M_{\sol}} f_{\sol}^{j-\mu}|(\nabla^{g_{\sol}})^j (\zeta_s^h)^{\ast}T|_{g_{\sol}}^2 \leq C_{\ell_1,\ell_2,m}(g_{\sol},\mu)\sum_{j=0}^{m} \sup_{M_{\sol}} f_{\sol}^{j-\mu}|(\nabla^{g_{\sol}})^jT|_{g_{\sol}}^2.
\end{equation*}

\item \label{modelmetrics3} 
for each $j\in\N$ and $h,k\in H$ with $\|h\|_H,\|k\|_H \leq 1$, 
\begin{equation*}
    \sup_{M_{\sol}}f_{\sol}^{\frac{j}{2}}|(\nabla^{g_{\sol}})^j(\omega_{h}-\omega_{k})|_{\omega_{\sol}}\leq C_j(g_{\sol})\|h-k\|_{H},
\end{equation*}

\item \label{modelmetrics4} there exists $C(g_{\sol})\in (1,\infty)$ such that for all $h,k\in H$ with $\|h\|_H,\|k\|_H \leq 1$,
\begin{equation*}
    |(u_{h}-u_{k})-(h-k)| \leq C(g_{\sol})(\|h\|_H + \|k\|_H)\|h-k\|_H f_{\sol}.
\end{equation*}
\end{enumerate}
\end{prop}

\begin{proof} \ref{modelmetrics0}
For any $t\in [0,1]$, we have
\begin{equation*} \partial_t \omega_{th}=(\zeta_t^h)^{\ast}\mathcal{L}_{\frac{1}{2}\nabla^{g_{\sol}}h}\omega_{\sol}=(\zeta_t^h)^{\ast}\sqrt{-1}\partial \overline{\partial}h .\end{equation*}
By Lemma \ref{lem:eigenfunctiongrowth} \ref{lem:eigenfunctiongrowth2}, we have
\begin{equation*}
     -C(g_{\sol})\|h\|_H\omega_{\sol}\le \sqrt{-1}\partial\bar\partial h\le C(g_{\sol})\|h\|_H\omega_{\sol}
\end{equation*}
for all $h\in H$, so combining expressions yields 
\begin{equation*}
   \partial_t (e^{-C(g_{\sol})\|h\|_Ht}\omega_{th}) \leq 0.
\end{equation*}
Upon integrating over $t\in [0,1]$, this yields $\omega_h \leq e^{C(g_{\sol})\|h\|_H}\omega_{\sol}$. After pulling back by $\zeta_1^h$, we obtain the remaining inequality.  

\ref{modelmetrics1}:
Recall that $u_h=\int_0^1(\zeta^h_s)^*hds$, where $(\zeta_s^h)_{s\in \mathbb{R}}$ is the flow of $\frac{1}{2}\nabla^{g_{\sol}}h$. Now we compute 
\begin{equation*}
   \begin{split}
        |u_h-u_k|&=\left|\int_0^1\left((\zeta^h_s)^*h-(\zeta^h_s)^*k+(\zeta^h_s)^*k-(\zeta^k_s)^*k\right)ds\right|\\
        &\le \int_0^1|(\zeta^h_s)^*(h-k)|ds+\int_0^1|(\zeta^h_s)^*k-(\zeta^k_s)^*k|ds\\
   \end{split}
\end{equation*}
By Lemma \ref{lem:eigenfunctiongrowth} and Lemma \ref{lem; equiv of potentials} \ref{lem; equiv of potentials1}, there is a uniform constant $C(g_{\sol})\in (1,\infty)$ such that for all $h\in H$ with $\|h\|_H \leq 1$, we have
\begin{equation*}
    |h|+\sqrt{f_{\sol}}|\nabla^{g_{\sol}}h|_{g_{\sol}}\le C\|h\|_{H}f_{\sol},\qquad \frac{1}{C(g_{\sol})}f_{\sol}\le f_h\le C(g_{\sol})f_{\sol}.
\end{equation*}
Hence, for all $t\in[0,1]$, we have
\begin{equation*}
    (\zeta^h_t)^*|h-k|\le C(g_{\sol})\|h-k\|_{H}f_{th}\le C(g_{\sol})\|h-k\|_H f_{\sol}.
\end{equation*}
To estimate the other term, we use that
\begin{equation*}
        (\zeta_s^h)^*k-(\zeta_s^k)^*k=\int_0^1\frac{\partial}{\partial \tau}\left((\zeta^{k+\tau(h-k)}_s)^*k\right)d\tau = \int_0^1 (\zeta^{k+\tau(h-k)}_s)^*(Z^\tau_s\cdot k)d\tau,
\end{equation*}
where $Z_s^{\tau}:=\frac{\partial}{\partial \tau}\zeta^{k+\tau(h-k)}_s\circ\zeta^{k+\tau(h-k)}_{-s}$. By Lemma \ref{tech lemma on lie derivative}, we have
\begin{equation} \label{eq:variation vector field model metrics}
    Z_s^\tau=\frac{1}{2}\int_0^s(\zeta_{s-\sigma}^{k+\tau(h-k)})_*(\nabla^{g_{\sol}}(h-k))d\sigma .
\end{equation}
 As a consequence of \ref{modelmetrics0}, for any tensor $T$ on $M_{\sol}$, and any $h\in H$, we can estimate
\begin{equation} \label{eq:usefuldoodad}
    |(\zeta_t^h)^{\ast}T|_{g_{\sol}}\leq C(g_{\sol}) |(\zeta_t^h)^{\ast}T|_{(\zeta_t^h)^{\ast}g_{\sol}}=C(g_{\sol})(\zeta_t^h)^{\ast}|T|_{g_{\sol}}.
\end{equation}
If $\|h\|_H,\|k\|_H \leq 1$, we thus have
\begin{align*}
    |(\zeta_{s-\sigma}^{k+\tau(h-k)})_{\ast}(\nabla^{g_{\sol}}(h-k))|_{g_{\sol}} (x) &\le C(g_{\sol})|\nabla^{g_{\sol}}(h-k)|_{g_{\sol}}(\zeta_{\sigma-s}^{k+\tau(h-k)}(x))\\
    & \leq C(g_{\sol})\|h-k\|_H \sqrt{f_{\sol}},
\end{align*}
where, again, we used Lemma \ref{lem:eigenfunctiongrowth} and Lemma \ref{lem; equiv of potentials} \ref{lem; equiv of potentials1} for the second inequality. Integrating yields

\begin{equation*}
    |Z_s^\tau|_{g_{\sol}}\leq 
    \frac{1}{2}
    Cs\|h-k\|_H \sqrt{f_{\sol}}. 
\end{equation*}
Therefore, for all $s\in[0,1]$ we have
\begin{equation*}
     |(\zeta_s^h)^*k-(\zeta_s^k)^*k|\le Cs\|k\|_H\|h-k\|_H f_{\sol}.
\end{equation*}
Hence, if $\|h\|_H,\|k\|_H\le 1$ we conclude that
\begin{equation*}
    |u_h-u_k|\le C\|h-k\|_H f_{\sol}.
\end{equation*}

\ref{modelmetrics2.5} We proceed by induction, with the case $m=0$ immediate from \eqref{eq:usefuldoodad}. Suppose now that $m\in\N^\times$ and that the claim holds with
$m$ replaced by $m-1$. Note that the same estimate holds for $s\in [-1,1]$ by replacing $h$ with $-h$. By homogeneity, we may assume that
\begin{equation}\label{eq:normalized weighted tensor norm}
\sum_{j=0}^{m}
\sup_{M_{\sol}}
f_{\sol}^{j-\mu}
\left|
    \bigl(\nabla^{g_{\sol}}\bigr)^jT
\right|_{g_{\sol}}^2
=1.
\end{equation}

Fix $x\in M_{\sol}$, set $x_s := \zeta_{-s}^h(x)$, $V:=\frac{1}{2}\nabla^{g_{\sol}}h$, $T_s:= (\zeta_s^h)^{\ast}T$, and define
\begin{equation*}\eta(s):=
\left|(\nabla^{g_{\sol}})^mT_s\right|_{g_{\sol}}^2(x_s)\end{equation*}
for all $s\in [0,1]$. 
Since $\partial_sT_s=\mathcal{L}_VT_s$, we compute
\begin{align*}
\eta'(s)
={}&
2\left\langle
    (\nabla^{g_{\sol}})^m\mathcal{L}_VT_s,
     (\nabla^{g_{\sol}})^mT_s
\right\rangle_{g_{\sol}}(x_s)
-
\left(
    V\cdot \left| (\nabla^{g_{\sol}})^mT_s\right|_{g_{\sol}}^2
\right)(x_s)
\\
={}&
2\left\langle
     (\nabla^{g_{\sol}})^m\mathcal{L}_V T_s
    -
     \nabla^{g_{\sol}}_V  (\nabla^{g_{\sol}})^mT_s,
     (\nabla^{g_{\sol}})^mT_s
\right\rangle_{g_{\sol}}(x_s).
\end{align*}
As usual, let $T'\ast T''$ denote a linear combination of
contractions of $T'$ and $T''$ with $g_{\sol}$ and
$g_{\sol}^{-1}$. Commuting covariant derivatives gives
\begin{align} (\nabla^{g_{\sol}})^m\mathcal{L}_VT_s
={}&
\nabla^{g_{\sol}}_V (\nabla^{g_{\sol}})^mT_s
+
\sum_{j=0}^{m}
(\nabla^{g_{\sol}})^{j+2}h
\ast
(\nabla^{g_{\sol}})^{m-j}T_s
\nonumber\\
&+
\sum_{j=0}^{m-1}
(\nabla^{g_{\sol}})^j
\left(
    \nabla^{g_{\sol}} h
    \ast
    \Rm(g_{\sol})
    \ast
    (\nabla^{g_{\sol}})^{m-1-j}T_s
\right),
\label{eq:commuted lie derivative}
\end{align}
so that $\xi(s):=\sqrt{\eta(s)}$ is a Lipschitz function satisfying the following for a.e. $s \in [0,1]$:
\begin{equation}
\label{eq:eta derivative expanded}
\begin{aligned}
|\xi'(s)|
\leq{}&
C_m
\sum_{j=0}^{m}
\left|(\nabla^{g_{\sol}})^{j+2}h\right|_{g_{\sol}}(x_s)
\left|(\nabla^{g_{\sol}})^{m-j}T_s\right|_{g_{\sol}}(x_s)
\\
&+
C_m
\sum_{j=0}^{m-1}
\sum_{\substack{a,b,c\geq 0\\a+b+c=j}}
\left|(\nabla^{g_{\sol}})^{a+1}h\right|_{g_{\sol}}(x_s)
\left|(\nabla^{g_{\sol}})^b\Rm(g_{\sol})\right|_{g_{\sol}}(x_s)
\left|
    (\nabla^{g_{\sol}})^{m-1-j+c}T_s
\right|_{g_{\sol}}(x_s).
\end{aligned}
\end{equation}
By Lemma
\ref{lem:eigenfunctiongrowth} \ref{lem:eigenfunctiongrowth2}, Lemma \ref{lem:basicACfacts} \ref{lem:basicACfacts3} and $\|h\|_H\leq 1$, we have
\begin{align*}
\left|(\nabla^{g_{\sol}})^{j+2}h\right|_{g_{\sol}}
\leq
C_j(g_{\sol})f_{\sol}^{-\frac{j}{2}}, \quad 
\left|(\nabla^{g_{\sol}})^{a+1}h\right|_{g_{\sol}}
\leq
C_a(g_{\sol})f_{\sol}^{\frac{1-a}{2}}, \quad
\left|(\nabla^{g_{\sol}})^b\Rm(g_{\sol})\right|_{g_{\sol}}
\leq
C_b(g_{\sol})f_{\sol}^{-\frac{b+2}{2}}.
\end{align*}
Moreover, the induction hypothesis and
\eqref{eq:normalized weighted tensor norm} imply that
\begin{equation}\label{eq:lower derivative pullback bound}
\left|(\nabla^{g_{\sol}})^rT_s\right|_{g_{\sol}}(x_s)
\leq
C_r(g_{\sol})
f_{\sol}(x_s)^{\frac{\mu-r}{2}}
\end{equation}
for every $r\in \{0,...,m-1\}$. The $j=0$ term in the first sum in
\eqref{eq:eta derivative expanded} is bounded by $C_m\xi(s)$, so combining expressions yields
\begin{equation*}
|\xi'(s)|
\leq
C_m(g_{\sol})
\left( \xi(s) + f_{\sol}^{\frac{\mu-m}{2}}(x_s)
\right) = C_m(g_{\sol})\left( \xi(s)+f_{-sh}^{\frac{\mu-m}{2}}(x) \right).
\end{equation*}
Integrating over $s\in [0,1]$ and using Lemma
\ref{lem; equiv of potentials} \ref{lem; equiv of potentials1}
gives
\begin{equation*}
\left|(\nabla^{g_{\sol}})^mT_s\right|_{g_{\sol}}(x_s)
\leq
C_m(g_{\sol})
f_{\sol}^{\frac{\mu-m}{2}}(x).
\end{equation*}
Since $x\in M_{\sol}$ was arbitrary, another application of Lemma \ref{lem; equiv of potentials}  \ref{lem; equiv of potentials1} proves the estimate
at order $m$, so the claim follows by induction.

\ref{modelmetrics3} For $t\in[0,1]$, set $h_t := k+t(h-k)$, and let $Z_s^t$ be as in the proof of \ref{modelmetrics1}. 
Lemma
\ref{lem:eigenfunctiongrowth}\ref{lem:eigenfunctiongrowth2}
then gives
\begin{equation}\label{eq:weighted derivatives gradient h minus k}
\sum_{j=0}^{m+1}
\sup_{M_{\sol}}
f_{\sol}^{j-1}
\left|
    \bigl(\nabla^{g_{\sol}}\bigr)^j
    \nabla^{g_{\sol}}(h-k)
\right|_{g_{\sol}}^2
\leq
C_{m+1}(g_{\sol})
\|h-k\|_H^2.
\end{equation}
Applying \ref{modelmetrics2.5} with
$\mu=1$, and then using Jensen's inequality in
\eqref{eq:variation vector field model metrics} yields
\begin{equation}\label{eq:weighted lie derivative estimate}
\sum_{j=0}^{m}
\sup_{M_{\sol}}
f_{\sol}^{j}
\left|
    \bigl(\nabla^{g_{\sol}}\bigr)^j
    \mathcal{L}_{Z_1^t}\omega_{\sol}
\right|_{g_{\sol}}^2\leq \sum_{j=0}^{m+1}
\sup_{M_{\sol}}
f_{\sol}^{j-1}
\left|
    \bigl(\nabla^{g_{\sol}}\bigr)^jZ_1^t
\right|_{g_{\sol}}^2
\leq
C_{m+1}(g_{\sol})
\|h-k\|_H^2,
\end{equation}
uniformly for $t\in[0,1]$. 
Applying \ref{modelmetrics2.5} with
$\mu=0$ and using
\eqref{eq:weighted lie derivative estimate}, we obtain
\begin{align*}
f_{\sol}^{\frac{m}{2}}
\left|
    \frac{\partial}{\partial t}
    \bigl(\nabla^{g_{\sol}}\bigr)^m\omega_{h_t}
\right|_{g_{\sol}}
=
f_{\sol}^{\frac{m}{2}}
\left|
    \bigl(\nabla^{g_{\sol}}\bigr)^m
    \left(\zeta_1^{h_t}\right)^*
    \mathcal{L}_{Z_1^t}\omega_{\sol}
\right|_{g_{\sol}}\leq
C_m(g_{\sol})\|h-k\|_H,
\end{align*}
so the claim follows by integrating over $t\in[0,1]$.

\ref{modelmetrics4} For $h\in H$, we compute
\begin{equation*}
    u_h-h=\int_0^1((\zeta_t^h)^*h-h)dt=\int_0^1\int_0^t(\zeta_s^h)^*|\partial  h|_{g_{\sol}}^2dsdt,
\end{equation*}
so that for any $h,k\in H$, 
\begin{equation*}
  \begin{split}
        (u_h-h)-(u_k-k)&=\int_0^1\int_0^t\left((\zeta_s^h)^*|\partial h|_{g_{\sol}}^2-(\zeta_s^k)^*|\partial k|_{g_{\sol}}^2\right)dsdt\\
        &=\int_0^1\int_0^t\left((\zeta_s^h)^*(|\partial h|_{g_{\sol}}^2-|\partial k|_{g_{\sol}}^2)+(\zeta_s^h)^*|\partial k|_{g_{\sol}}^2-(\zeta_s^k)^*|\partial k|_{g_{\sol}}^2\right)dsdt
  \end{split}
\end{equation*}
For the first term, we have
\begin{equation*}
\begin{split}
     \left|  |\partial h|_{g_{\sol}}^2-|\partial k|^2_{g_{\sol}}\right|&\le ( |\partial h|_{g_{\sol}}+ |\partial k|_{g_{\sol}})|\partial(h-k)|_{g_{\sol}}\\
     &\le C(\|h\|_H+\|k\|_H)\|h-k\|_H f_{\sol}
\end{split}
\end{equation*}
which implies the following for $\|h\|_H \leq 1$:
\begin{equation*}
    \int_0^1\int_0^t|(\zeta_s^h)^*|\partial h|_{g_{\sol}}^2-(\zeta_s^h)^*|\partial k|_{g_{\sol}}^2|dsdt\le C(\|h\|_H+\|k\|_H)\|h-k\|_H f_{\sol}.
\end{equation*}
For the remaining term, we use that if $\|k+r(h-k)\|_H \leq 1$ for all $r\in[0,1]$, then
\begin{equation*}
    \begin{split}
       \left| (\zeta_s^h)^*|\partial k|_{g_{\sol}}^2-(\zeta_s^k)^*|\partial k|_{g_{\sol}}^2\right|&=\left|\int_0^1\frac{\partial}{\partial r}\left((\zeta_s^{k+r(h-k)})^*|\partial k|_{g_{\sol}}^2\right)dr\right|\\
        &=\left|\int_0^1(\zeta_s^{k+r(h-k)})^*(Z_s^r\cdot|\partial k|_{g_{\sol}}^2)dr\right|\\
        &\le Cs \|h-k\|_H\|k\|_H^2 f_{\sol}.    
    \end{split}
\end{equation*}
If $h,k\in H$ satisfy $\|h\|_H,\|k\|_H \leq 1$, it follows that
\begin{equation*}
      |(u_h-h)-(u_k-k)|\le C\|h-k\|_H(\|h\|_H+\|k\|_H)f_{\sol}.
\end{equation*}
This completes the proof of \ref{modelmetrics4}.
\end{proof}

\end{document}
